% denpy manual: data formats, zarr helpers and the TIFF stack conversion.
% Compile with:  tectonic denpy_manual.tex
\documentclass[11pt,a4paper]{article}

\usepackage{fontspec}
\usepackage[margin=2.3cm]{geometry}
\usepackage{xcolor}
\usepackage[most]{tcolorbox}
\usepackage{listings}
\usepackage{booktabs,tabularx,longtable,array}
\usepackage{enumitem}
\usepackage{titlesec}
\usepackage{fancyhdr}
\usepackage[hidelinks]{hyperref}

\definecolor{accent}{HTML}{1F5A96}
\definecolor{termbg}{HTML}{F4F6F8}
\definecolor{termframe}{HTML}{B8C2CC}
\definecolor{notebg}{HTML}{EAF2FB}
\definecolor{warnbg}{HTML}{FDF1E4}
\definecolor{warnframe}{HTML}{D9822B}

\hypersetup{colorlinks=true,linkcolor=accent,urlcolor=accent}
\setlist{itemsep=2pt,topsep=4pt}
\setlength{\parindent}{0pt}
\setlength{\parskip}{5pt}
\setlength{\emergencystretch}{3em}

\titleformat{\section}{\Large\bfseries\sffamily\color{accent}}{\thesection}{0.8em}{}
\titleformat{\subsection}{\large\bfseries\sffamily}{\thesubsection}{0.8em}{}
\titleformat{\subsubsection}{\normalsize\bfseries\sffamily}{\thesubsubsection}{0.8em}{}

\pagestyle{fancy}
\fancyhf{}
\fancyhead[L]{\small\sffamily denpy --- Manual}
\fancyhead[R]{\small\sffamily\nouppercase{\leftmark}}
\fancyfoot[C]{\small\thepage}
\renewcommand{\headrulewidth}{0.4pt}

\lstset{
  basicstyle=\ttfamily\small,
  breaklines=true,
  breakatwhitespace=false,
  postbreak=\mbox{\textcolor{gray}{$\hookrightarrow$}\space},
  columns=fullflexible,
  keepspaces=true,
  showstringspaces=false,
  tabsize=4,
  upquote=true,
  language={},
}

\tcbset{terminalstyle/.style={listing only, colback=termbg, colframe=termframe, boxrule=0.4pt,
  arc=2pt, left=4pt, right=4pt, top=2pt, bottom=2pt, breakable}}
\newtcblisting{terminal}{terminalstyle}
\newtcblisting{terminalsmall}{terminalstyle, listing options={basicstyle=\ttfamily\footnotesize}}
\newtcolorbox{note}{colback=notebg, colframe=accent, boxrule=0.4pt, arc=2pt, left=6pt, right=6pt, breakable}
\newtcolorbox{warning}{colback=warnbg, colframe=warnframe, boxrule=0.4pt, arc=2pt, left=6pt, right=6pt, breakable}

\newcommand{\code}[1]{\texttt{\detokenize{#1}}}

\begin{document}
\sloppy

\begin{titlepage}
\centering
\vspace*{4cm}
{\Huge\bfseries\sffamily\color{accent} denpy\par}
\vspace{0.6cm}
{\Large\sffamily Data formats, zarr storage and TIFF stack conversion\par}
\vspace{1.2cm}
{\large Manual for denpy 3.0.0\par}
\vfill
{\small Vojtěch Kulvait, Helmholtz-Zentrum Hereon\par}
{\small November 2026\par}
\end{titlepage}

\tableofcontents
\newpage

% ====================================================================
\section{Overview}

denpy is the Python package for the data formats of the KCT tomography tools: the DEN format of
the KCT C++ tools, zarr stores (directory or zip), and the conversion of stacks of TIFF files into
either. It has no knowledge of beamlines or instruments: the parsing of scan metadata (HDF5 files of
PETRA III, nanoCT logs) lives in lautools, which builds the conversion specifications described here
(see the lautools API manual).

\begin{longtable}{@{}p{3.2cm}p{11.6cm}@{}}
\toprule
\textbf{Module} & \textbf{Content}\\
\midrule
\endhead
\code{DEN} & DEN files: headers, frames, whole arrays, files of more than 3 dimensions (\S\ref{sec:den}).\\
\code{ZAR} & zarr codecs: \code{get_compressor}, \code{codec_arguments}, \code{register_codecs} (\S\ref{sec:zar}).\\
\code{ZARZS} & Zip store that threads write and read in parallel, safe appends (\S\ref{sec:zarzs}).\\
\code{ZARPYRAMID} & Multiscale pyramids with OME-NGFF metadata (\S\ref{sec:pyramid}).\\
\code{TIFF} & TIFF stacks to zarr or DEN from a JSON specification, with virtual stacks (lambdas, reductions) (\S\ref{sec:tiff}).\\
\code{PETRA} & Legacy parser of PETRA III scan HDF5 files; the object oriented implementation is \code{lautools.PETRA} (\S\ref{sec:petra}).\\
\code{COR}, \code{ECC} & Rotation centre detection on sinograms; image registration (OpenCV ECC).\\
\code{PHANTOM} & Shepp--Logan like phantoms in 2D and 3D, stored as DEN.\\
\code{DICOM}, \code{UTILS} & DICOM helpers; parameter files.\\
\bottomrule
\end{longtable}

Command line: \code{convertTiffStack} (\S\ref{sec:cli}).

\subsection{Installation}

\begin{terminal}
pip install denpy              # DEN, PETRA (h5py, pandas), phantoms, ...
pip install "denpy[zarr]"      # + zarr, tifffile, imagecodecs: ZAR, ZARZS, TIFF, convertTiffStack
\end{terminal}

On Maxwell the shared Lau environment has denpy with all extras. A reader of zarr arrays written with
image codecs (htj2k, jpeg2k, jpegxl, ...) must register the codecs first:
\code{ZAR.register_codecs()}; otherwise zarr reports e.g.\ \emph{Unknown codec: imagecodecs\_htj2k}.

\subsection{Tests}

Every module has a unit test next to it (\code{denpy/test_*.py}):

\begin{terminal}
python -m pytest denpy/          # in the repository root
\end{terminal}

% ====================================================================
\section{DEN files}\label{sec:den}

A DEN file is a 4096 byte header followed by the raw array. The extended header holds the number of
dimensions, the element size, the majority (x or y major) and the element type, followed by the
dimensions \emph{dimspec} in the order $(X, Y, Z_1, Z_2, \dots)$: the first dimension changes fastest.
For a numpy array of shape \code{(..., Y, X)} the dimspec is the shape reversed. Supported types are
uint8, uint16, int16, uint32, int32, uint64, int64, float32 and float64. Legacy files have a 6 byte
header with three uint16 dimensions.

\begin{longtable}{@{}p{6.2cm}p{8.6cm}@{}}
\toprule
\textbf{Function} & \textbf{Purpose}\\
\midrule
\endhead
\code{readHeader(f)} & Dict: \code{dimspec}, \code{shape} (numpy order), \code{type}, \code{offset}, \code{dimcount}, \code{majority}.\\
\code{getNumpyArray(f)} & The whole array.\\
\code{getFrame(f, k, ...)} & Frame $k$ (for more than 3 dimensions a list of indices in dimspec order), optionally a region and a dtype.\\
\code{writeEmptyDEN(f, dimspec, elementtype)} & A file of zeros of the given size (sparse).\\
\code{writeFrame(f, k, data, force)} & Write frame $k$ into an existing file.\\
\code{storeNdarrayAsDEN(f, array)} & Write a whole numpy array.\\
\code{npDtypeToDenDataType}, \code{denDataTypeToNpDtype} & Type mapping.\\
\bottomrule
\end{longtable}

\begin{terminal}
from denpy import DEN
info = DEN.readHeader("img.den")          # info["shape"] == (frames, Y, X)
frame = DEN.getFrame("img.den", 10)
DEN.storeNdarrayAsDEN("out.den", array, force=True)
\end{terminal}

% ====================================================================
\section{Zarr helpers}

\subsection{Codecs: ZAR}\label{sec:zar}

\code{ZAR.get_compressor(name, clevel, zarrv2=False, dtype=None, **kwargs)} returns the codec for a
name: \code{none}, \code{zstd}, \code{lz4}, \code{gzip}, \code{blosc-*} (zstd, lz4, lz4hc, zlib,
snappy, blosclz; with bit shuffle), and the image codecs \code{htj2k}, \code{jpeg2k}, \code{jpegxl},
\code{jpegxr}, \code{avif}, \code{zfp}, \code{sz3}. For zarr v3 it returns a codec list; pass it
through \code{ZAR.codec_arguments(codec)} to \code{create_array}, which makes an image codec the
serializer and the others compressors.

\begin{note}
\textbf{Levels.} For htj2k, jpeg2k and jpegxl \code{clevel 0} is lossless. For htj2k a level
$>0$ is a quantization step relative to the full range: 0.001 is already very lossy. Blosc with level
0 does not compress. Bit shuffle is 1.7--4 times faster than byte shuffle for detector data.
Measured on 5120$\times$3840 uint16 frames, lossless htj2k is the best compromise of size and speed.
\end{note}

\subsection{Parallel zip store: ZARZS}\label{sec:zarzs}

\code{ZARZS.ZipStore} is a drop-in for \code{zarr.storage.ZipStore}. zarr's own zip store writes and
reads under one lock, so a zip is much slower than a directory store. In ZARZS each thread computes
the CRC of its chunk, reserves the space under the lock (microseconds) and writes with \code{pwrite}
outside it; reads use \code{pread}. The result is a standard zip (entries stored uncompressed, the
chunks are compressed by zarr) that zarr, Python's zipfile and \code{unzip} read.

\begin{terminal}
from denpy import ZARZS
store = ZARZS.ZipStore("volume.zar", mode="w")   # "r", "w" or "a"
group = zarr.open_group(store, mode="w")
...                                              # write from as many threads as useful
store.close()                                    # writes the directory of the zip
\end{terminal}

Beyond zarr's zip store:
\begin{itemize}
\item a key written again replaces its entry (no duplicate names in the zip), \code{delete} and
  \code{delete_dir} work, so arrays can be overwritten;
\item \textbf{safe append}: in mode \code{"a"} new entries go after the end of the file and the
  original size is recorded in \code{<zip>.zarzs-append} until \code{close()}. If the process is
  killed in between, the store refuses to open the zip and \code{ZARZS.recover(path)} truncates it to
  the state before the append, byte for byte;
\item a store that is garbage collected without \code{close()} still leaves a valid zip.
\end{itemize}
Limits: threads of one process only (a writable store can not be pickled); entries are uncompressed;
compressed entries of zips written by other tools are read through zipfile.

\subsection{Pyramids: ZARPYRAMID}\label{sec:pyramid}

\code{ZARPYRAMID.write_pyramid(output, frames, shape, dtype, spec)} writes a Z,Y,X stack frame by frame
as a multiscale pyramid (arrays \code{0}, \code{1}, \dots), by default with OME-NGFF metadata (version
0.5 for zarr v3, 0.4 for v2). The \code{PyramidSpec} sets \code{levels} (5), \code{factor} (2),
\code{axes} (\code{yx}: Z kept, or \code{zyx}) and \code{method} (\code{mean}, \code{nearest},
\code{max}, \code{min}). \code{frames} is an iterable or a callable $k \mapsto$ frame. An output path
ending with \code{.zip} is written directly with ZARZS; \code{PyramidWriter} writes level by level
for callers that produce the frames themselves.

% ====================================================================
\section{TIFF stacks to zarr or DEN}\label{sec:tiff}

\code{denpy.TIFF} converts stacks of TIFF files into arrays of a zarr store or DEN files as described
by a JSON specification. The full reference of the format is \code{docs/TIFFSTACK_SPEC.md}; this
section explains it with examples.

\subsection{The three parts of a specification}

\begin{longtable}{@{}p{3.3cm}p{2.3cm}p{9.2cm}@{}}
\toprule
\textbf{Part} & \textbf{Layer} & \textbf{What it says}\\
\midrule
\endhead
\code{tiffstacks} & reading & Which files form a stack, their frame shape and dtype, which frames are known to be missing.\\
\code{virtualstacks} & computing & Stacks computed from other stacks: a lambda applied frame by frame, or a reduction to one frame; not written.\\
\code{arrays} / \code{array} & writing & What is written: a stack as it is or computed like a virtual stack, with the storage settings.\\
\code{output} & writing & Zarr (zip or directory store) or DEN, mode, attributes.\\
\bottomrule
\end{longtable}

The shape and dtype of every array follow from the specification alone, before any TIFF is read, so
a specification is validated completely up front: the files must exist, frame counts must match,
lambdas are checked and run once on frames of zeros. While writing, every frame is checked against
the declared shape and dtype. All ranges are half-open, \code{[start, stop)}, as in Python; arrays are
in numpy order (a frame is \code{[Y, X]}).

\subsection{A first example}

\begin{terminalsmall}
{
  "spec": "denpy.tiffstack", "version": 1,
  "base_dir": "/gpfs/.../raw/scan_042",
  "tiffstacks": {
    "dar": {"frames": ["dar_00000.tif", "dar_00001.tif", "dar_00002.tif"],
            "frame_shape": [3840, 5120], "dtype": "uint16"},
    "ref": {"frames": [{"pattern": "ref_{:05d}.tif", "range": [0, 20]},
                       {"pattern": "ref_{:05d}.tif", "range": [1520, 1540]}],
            "framecount": 40, "frame_shape": [3840, 5120], "dtype": "uint16"},
    "img": {"frames": [{"pattern": "img_{:05d}.tif", "range": [0, 1500]}], "framecount": 1500,
            "frame_shape": [3840, 5120], "dtype": "uint16"}
  },
  "arrays": {"dar": {"stack": "dar"}, "ref": {"stack": "ref"}, "img": {"stack": "img"}},
  "defaults": {"zarr": {"codec": {"name": "htj2k", "clevel": 0}}},
  "output": {"format": "zarr", "store": "zip", "path": "scan_042.zar"}
}
\end{terminalsmall}

\begin{terminal}
convertTiffStack scan_042.json --dry-run     # validate, print shapes and dtypes
convertTiffStack scan_042.json -j 32 --report report.json
\end{terminal}

\subsection{Tiff stacks}

\begin{longtable}{@{}p{3cm}p{11.8cm}@{}}
\toprule
\textbf{Key} & \textbf{Meaning}\\
\midrule
\endhead
\code{frames} & List of entries in stack order: \code{"a.tif"}; \code{{"file": "a.tif", "page": 3}} (a page of a multi-page TIFF); \code{null} (a frame that should exist but has no file); \code{{"pattern": "img_{:05d}.tif", "range": [start, stop(, step)]}}.\\
\code{framecount} & Number of frames; \textbf{required} with a pattern entry, checked when given.\\
\code{frame_shape}, \code{dtype} & Required: \code{[Y, X]} and dtype of every frame.\\
\code{fill_value} & Value of the \code{null} frames; required when there is one.\\
\bottomrule
\end{longtable}

Everything a stack lists must be there: a missing or unreadable file, or a frame of another shape or
dtype, is an error. The only frames not read from a file are the \code{null} entries the author of the
specification put there knowingly; the report lists them. Patterns are Python \code{str.format} with
exactly one integer field. Relative file names resolve against \code{base_dir}, by default the
directory of the specification file.

\subsection{Virtual stacks: lambdas and reductions}

A virtual stack computes frame $k$ from other stacks:

\begin{terminalsmall}
"virtualstacks": {
  "dar_median": {"stack": "dar", "reduce": "median", "dtype": "uint16"},
  "img_darksub": {"stack": "img", "dtype": "uint16", "clip": true,
                  "lambda": "lambda x, dark: x.astype('int32') - dark.astype('int32')",
                  "args": ["dar_median"]}
}
\end{terminalsmall}

\begin{itemize}
\item \code{reduce}: \code{mean}, \code{median}, \code{min}, \code{max}, \code{sum} of the frames: a
  single frame (2D). \code{median} holds the stack in memory; validation rejects it above
  \code{max_memory} (default a quarter of the RAM). Null frames are left out.
\item \code{lambda}: one Python lambda expression with $1+$\code{len(args)} parameters; \code{np} is
  available. It gets frame $k$ of \code{stack} and one value per \code{args} entry: a 2D source (a
  reduction) gives the same frame for every $k$, a source with the same number of frames its frame
  $k$, a number the number.
\item \code{function}: \code{"module:name"}, an importable function with the same signature, for
  anything longer than a line (in the Python API also a callable).
\item \code{dtype}: the result is converted; float to integer rounds to nearest; values outside the
  range of an integer dtype are an error, or clipped with \code{"clip": true}.
\end{itemize}

Examples of lambdas:
\begin{terminalsmall}
"lambda x: x[100:3940, :]"                                         crop
"lambda x: x.reshape(1920, 2, 2560, 2).mean(axis=(1, 3))"           2x2 binning
"lambda x, dark, flat: (x - dark) / np.maximum(flat - dark, 1)"     flat field correction
\end{terminalsmall}

\begin{warning}
A specification with lambdas or functions is code: run only specifications you trust, as any script.
\end{warning}

Names: tiff stacks and virtual stacks are the \emph{sources} and share one namespace; \code{stack}
and \code{args} refer to sources. Arrays are outputs, their names are paths in the output and may
repeat a source name (\code{"img": {"stack": "img"}}). To write a computed stack and use it elsewhere,
define it as a virtual stack and write it with \code{{"stack": name}}; it is computed once.

\subsection{Arrays and output}

An array is \code{{"stack": name}} or computed like a virtual stack, with the storage keys
\code{grid} (frames as more than one dimension: \code{[30, 50]} gives shape \code{(30, 50, Y, X)}),
\code{attributes}, \code{overwrite} and \code{zarr}: \code{codec} (\code{{"name": ..., "clevel": ...}},
the arguments of \code{ZAR.get_compressor}), \code{chunks} (default one frame per chunk; $-1$ is the
full extent), \code{shards}, \code{dimension_names}. \code{defaults} gives keys common to all arrays.
With \code{"array"} instead of \code{"arrays"} the zarr root is the array itself, or the DEN output a
single file.

\begin{longtable}{@{}p{2.6cm}p{12.2cm}@{}}
\toprule
\textbf{Output key} & \textbf{Meaning}\\
\midrule
\endhead
\code{format} & \code{"zarr"} or \code{"den"}.\\
\code{store} & Zarr: \code{"zip"} (ZARZS) or \code{"directory"}; the suffix of the path is not interpreted (\code{.zar} is recommended for zips).\\
\code{path} & The store; DEN: the file (\code{array}) or the directory of \code{<name>.den} files (\code{arrays}).\\
\code{mode} & \code{"w-"} (default) refuse existing output, \code{"w"} replace, \code{"a"} add arrays (\code{overwrite} for existing ones).\\
\code{zarr_format} & 3 (default) or 2.\\
\code{attributes} & Attributes of the root group.\\
\code{atomic} & Write to a temporary name and rename at the end (default true).\\
\code{provenance} & Store spec, expanded frame lists and report as attribute \code{denpy_tiffstack} of each array.\\
\bottomrule
\end{longtable}

\begin{note}
\textbf{A DEN directory is a container.} DEN files go into the directory, which may exist and hold
other files (e.g.\ a working directory): it is never removed. Each \code{<name>.den} is written under a
temporary name and renamed when all arrays are written; \code{"w-"} refuses existing files of the
arrays, \code{"w"} replaces just them. A zarr directory store with \code{"w"} is replaced as a whole.
\end{note}

\subsection{Python API}

\begin{terminal}
from denpy import TIFF
plan = TIFF.validate("scan.json")        # or a dict; SpecError if anything does not fit
print(plan.describe())
report = TIFF.execute(plan, threads=16)  # FrameError if a frame does not match
spec = TIFF.extend_spec(spec, "dark_correction.json")
spec = TIFF.spec_from_files(files, output)   # one stack, one array
\end{terminal}

\code{execute} reads, computes and writes frames in worker threads (default a quarter of the CPUs,
16 to 32): a thread reads the TIFF frames of frame $k$, applies the lambdas, converts and writes
(zarr compresses in the same thread); a zip is written in parallel through ZARZS. Reductions are
computed first, in dependency order. The report has, per array, the shape, dtype, null frames, TIFF
bytes read, stored bytes, ratio to the raw size and time.

\code{extend_spec(spec, extension)} adds \code{tiffstacks} (new names only), \code{virtualstacks}
(added or replaced), \code{arrays} (added or replaced; \code{null} removes an array) and merges
\code{defaults}; this is how a specification generated from metadata gets processing added.

\subsection{convertTiffStack}\label{sec:cli}

\begin{terminalsmall}
convertTiffStack spec.json [-j THREADS] [--extend EXT.json ...] [--report r.json] [--dry-run]
                 [--check-headers] [--max-memory GiB] [--write-spec full.json] [--verbose]
convertTiffStack --files img_*.tif --zarr img.zar [--store zip|directory] [--zarr-format 2|3]
                 [--compression htj2k --clevel 0] [--dtype float32 [--clip]] [--sort natural]
                 [--mode w-|w] [--write-spec spec.json] [--dry-run]
convertTiffStack --files img_*.tif --den img.den ...
\end{terminalsmall}

The second form builds a one-array specification from the file list (shell glob order, or natural
order with \code{--sort natural}); \code{--write-spec} saves it as a starting point.
\code{--check-headers} compares the TIFF header of every file with the specification before writing.
The exit status is 1 with an \code{ERROR:} line for an invalid specification or a frame that does not
match.

\subsection{Performance}

On a 128-CPU Maxwell node with 2500 frames of 5120$\times$3840 (91.5\,GB of TIFF files), reading
the TIFF files takes about 120\,s; lossless htj2k into a zip takes about 140\,s, about the time of a
directory store; more than 32 threads do not help. In the tests, 16 threads convert 120 frames of
2048$^2$ to an htj2k zip five times faster than one thread; a dark corrected stack costs about 7\,\%
more than the raw one.

% ====================================================================
\section{PETRA (legacy)}\label{sec:petra}

\code{denpy.PETRA} parses PETRA III scan HDF5 files: \code{scanDataset}, \code{imageDataset},
\code{getExperimentInfo}, \code{beamCurrentDataset}, \code{findInsertionEvents}. It is the legacy
implementation, kept for the existing callers. All its functions read through one opening of the
file and accept a file name or an open \code{h5py} file or group.

The object oriented implementation is \code{lautools.PETRA}. It parses HDF5 files of all beamline
layouts (P05 2021 and 2025, P07) and the \code{params} group of zarr stores into one canonical
\code{ScanDescription}, implements every analysis once on it, and writes the description to zarr and
to HDF5 in the 2025 layout. It has the same module functions, so a caller switches with the import:

\begin{terminal}
from lautools import PETRA     # instead of: from denpy import PETRA
\end{terminal}

Its API is described in the \emph{lautools API manual} (\code{manual/lautools_api_manual.pdf} of
lautools), together with the conversion of PETRA III scans to zarr and DEN built on \code{denpy.TIFF}
(\code{tiffScanDataToZarrForMicrotomography}). The lautools tests compare both implementations on
synthetic scans and real beamline files.

\end{document}
